% \section{Teoria elementare di prodotti, somme, equalizzatori, coequalizzatori}
% \Todo{}
% La descrizione del prodotto di due insiemi data in \ref{prod_in_Set} si presta a essere generalizzata a una generica categoria.
% \begin{definition}
% 	\Todo{Prodotto in \(\ctC\)}
% \end{definition}
% \begin{definition}
% 	\Todo{Coprodotto in \(\ctC\)}
% \end{definition}
% Si vede ora che le due definizioni sono duali in questo senso:
% \begin{lemma}
% 	prodotti in \(\ctC\) sono coprodotti in \(\ctC^\op\); coprodotti in \(\ctC\) sono prodotti in \(\ctC^\op\).
% \end{lemma}
% \begin{theorem}
% 	\Todo{Prodotto come oggetto terminale.}
% \end{theorem}
% \begin{corollary}
% 	\Todo{Unicità a meno di iso del prodotto.}
% \end{corollary}
% \color{red}
% \begin{definition}[Prodotti binari]
% 	Sia \(\ctC\) una categoria, e siano \(A\) e \(B\) due oggetti di \(\ctC\). Un prodotto di \(A\) e \(B\) è una spanna
% 	\[
% 		\begin{tikzcd}
% 			A&P\ar[l, "p_1"']\ar[r, "p_2"]&B
% 		\end{tikzcd}
% 	\]
% 	tale che, per ogni altra spanna
% 	\[\begin{tikzcd}
% 			A&Q\ar[l, "f"']\ar[r, "g"]&B
% 		\end{tikzcd}
% 	\]
% 	esista un'unica freccia
% 	\begin{tikzcd}
% 		h\colon Q\ar[r] &P
% 	\end{tikzcd}
% 	tale che \(p_1\cmp h=f\) e \(p_2\cmp h=g\).
% \end{definition}
% È utile seguire la definizione sui  diagrammi:\\[2ex]
% \begin{tikzpicture}
% 	\node [mybox] (box){%
% 		\begin{minipage}[t][19ex]{0.40\textwidth}
% 			\begin{center}
% 				\begin{tikzcd}
% 					&Q\ar[dl, "f"']\ar[dr, "g"]
% 					\\
% 					A&P\ar[l, "p_1"']\ar[r, "p_2"]&B
% 				\end{tikzcd}\\[2ex]
% 				\emph{per ogni spanna \((Q,f,g)\) \\su \(A\) e \(B\)}...
% 			\end{center}
% 		\end{minipage}
% 	};
% \end{tikzpicture}%
% \hfill
% \begin{tikzpicture}
% 	\node [mybox] (box){%
% 		\begin{minipage}[t][19ex]{0.40\textwidth}
% 			\begin{center}
% 				\begin{tikzcd}
% 					&Q\ar[dl, "f"']\ar[dr, "g"]\ar[d, "h", dashed]
% 					\\
% 					A&P\ar[l, "p_1"']\ar[r, "p_2"]&B
% 				\end{tikzcd}\\[2ex]
% 				...\emph{esiste un'unica freccia \(h\)\\ tale che \(p_1\cmp h=f\) e \(p_2\cmp h=g\)}
% 			\end{center}
% 		\end{minipage}
% 	};
% \end{tikzpicture}%
% %
% \\[2ex]
% dove la freccia \(h\) è tratteggiata perché la sua esistenza è conseguenza della definizione di prodotto, e i due i triangoli commutativi corrispondono alle due equazioni che concludono la definizione.

% \medskip
% Per riferirci al prodotto di due oggetti \(A\) e \(B\) diremo che la tripla \((P,p_1,p_2)\) è una \emph{spanna universale} su  \(A\) e \(B\). Inoltre, chiamiamo la proprietà che definisce il prodotto di due oggetti \emph{proprietà universale del prodotto}.

% \begin{example}\label{esempio_spanna_universale_in_Set}
% 	Consideriamo gli insiemi
% 	\[
% 		A=\{0,1\}\qquad B=\{0,1,2\}\qquad P=\{0,1,2,3,4,5\}
% 	\]
% 	e le funzioni
% 	\[
% 		p_1\colon P\to A\qquad n\mapsto n\pmod 2
% 	\]
% 	\[
% 		p_2\colon P\to B\qquad n\mapsto n\pmod 3
% 	\]
% 	La spanna \((P, p_1, p_2)\) è universale, ovvero definisce un prodotto di \(A\) e \(B\).


% 	Dimostriamolo. Sia \(Q\) un insieme e \(f\colon Q\to A\), \(g\colon Q\to B\) due funzioni. Per ogni elemento \(q\in Q\) possiamo due elementi \(f(q)\in A\) e \(g(q)\in B\). Definiamo allora \(h(q)\in P\) come l'unico intero positivo di \(P\) tale che
% 	\[
% 		h(q)\equiv p_1(q) \pmod 2\qquad\text{e}\qquad h(q)\equiv p_2(q) \pmod 3
% 	\]
% 	Il fatto che tale intero esista e sia unico è conseguenza del ben noto teorema \emph{teorema cinese dei resti}\footnote{Si veda su qualunque manuale di algebra o matematica discreta.} Nel caso in esame, può valere la pena verificarlo direttamente sulla tabella qui sotto
% 	\begin{center}
% 		\begin{tabular}{|l|c|c|c|c|c|c|c|c|}
% 			\hline
% 			\(n\)         & \(0\) & \(1\) & \(2\) & \(3\) & \(4\) & \(5\) & \(h(q)\)   \\
% 			\hline
% 			\(n \pmod 2\) & \(0\) & \(1\) & \(0\) & \(1\) & \(0\) & \(1\) & \(p_1(q)\) \\
% 			\hline
% 			\(n \pmod 3\) & \(0\) & \(1\) & \(2\) & \(0\) & \(1\) & \(2\) & \(p_2(q)\) \\
% 			\hline
% 		\end{tabular}
% 	\end{center}
% 	e la proprietà universale del prodotto è dimostrata per la spanna \((P, p_1, p_2)\).
% \end{example}
% L'esempio precedente mostra come il prodotto di due insiemi \(A\) e \(B\) possa essere presentato in diversi modi.

% Si sarebbe potuto descrivere, come abbiamo fatto  in \ref{prod_in_Set},
% \[
% 	A\times B=\{(0,0),\,(0,1),\,(0,2),(1,0),\,(1,1),\,(1,2)\}
% \]
% insieme alle funzioni
% \[
% 	\pi_A\colon A\times B\to A\qquad (n,m)\mapsto n
% \]
% \[
% 	\pi_B \colon A\times B\to B\qquad (n,m)\mapsto m
% \]
% e questa costruzione è chiamata talvolta canonica\footnote{Sul termine \emph{canonico} non ci dilunghiamo; basti osservare che esso è di solito utilizzato in modo informale per intendere la costruzione più semplice o più utilizzata, non essendoci una definizione \emph{canonica} di canonico!}, così come canoniche sono dette le proiezioni \(\pi_A\) e \(\pi_B\).

% Diversamente, si può considerare la spanna universale \((P,p_1,p_2)\), come abbiamo fatto in \ref{esempio_spanna_universale_in_Set}. Ora, la tabella qui sopra può essere interpretata come una funzione biettiva
% \[
% 	\varphi\colon P\to A\times B \qquad n\mapsto (n\pmod2,n\pmod3)
% \]
% che commuta con le proiezioni, cioè  tale che \(\pi_A\cmp\varphi=p_1\) e \(\pi_B\cmp\varphi=p_2\).

% Questa è una proprietà generale.

% \begin{proposition}[Il prodotto è unico a meno di un unico isomorfismo]
% 	In una categoria \(\ctC\), date due spanne universali \((P,p_1,p_1)\) e \((Q,q_1,q_2)\) che insistano sulla stessa coppia di oggetti \(A,B\), esiste un unico isomorfismo \(h\colon P\to Q\) tale che \(q_1\cmp h=p_1\) e \(q_2\cmp h=p_2\).
% \end{proposition}
% Questa proprietà non è specifica dei prodotti binari, ma riguarda tutti i limiti e tutti i colimiti, e può essere facilmente ottenuta come caso particolare della \ref{???}. Riteniamo utile esplicitarne la dimostrazione, a fini didattici.
% \begin{proof}
% 	Poiché la spanna \((P,p_1,p_1)\) è universale, esiste un'unica freccia  \(h\colon Q\to P\) tale che \(p_1\cmp h=q_1\) e \(p_2\cmp h=q_2\). Dimostriamo che \(h\) è un isomorfismo esibendo un suo inverso. Per trovarlo ripetiamo il ragionamento precedente, questa volta utilizzando la spanna universale \((Q,q_1,q_2)\), e ottenendo un'unica freccia \(k\colon P\to Q\) tale che \(q_1\cmp =p_1\) e \(q_2\cmp h=p_2\). Sostituendo le nelle relazioni trovate (o \emph{impastando} i diagrammi in basso a sinistra) si ottiene che:\\[2ex]
% 	\begin{tikzpicture}
% 		\node [mybox] (box){%
% 			\begin{minipage}[t][28ex]{0.40\textwidth}
% 				\begin{center}
% 					\begin{tikzcd}
% 						&P\ar[ddl, "p_1"', bend right]\ar[ddr, "p_2", bend left]\ar[d, "k", dashed]
% 						\\
% 						&Q\ar[dl, "q_1"']\ar[dr, "q_2"]\ar[d, "h", dashed]
% 						\\
% 						A&P\ar[l, "p_1"']\ar[r, "p_2"]&B
% 					\end{tikzcd}\\[2ex]
% 					\(h\cmp k\) soddisfa le relazioni\\
% 					\(p_1\cmp h\cmp k = q_1\cmp k=p_1\)\\
% 					\(p_2\cmp h\cmp k = q_2\cmp k=p_2\)
% 				\end{center}
% 			\end{minipage}
% 		};
% 	\end{tikzpicture}%
% 	\hfill
% 	\begin{tikzpicture}
% 		\node [mybox] (box){%
% 			\begin{minipage}[t][28ex]{0.40\textwidth}
% 				\begin{center}
% 					\begin{tikzcd}
% 						&P\ar[ddl, "p_1"', bend right]\ar[ddr, "p_2", bend left]\ar[dd, "\id_P",]
% 						\\
% 						\\
% 						A&P\ar[l, "p_1"']\ar[r, "p_2"]&B
% 					\end{tikzcd}\\[4ex]
% 					...ma anche \(\id_P\) è tale che \\
% 					\(p_1\cmp \id_P =p_1\)\\
% 					\(p_2\cmp \id_P =p_2\)
% 				\end{center}
% 			\end{minipage}
% 		};
% 	\end{tikzpicture}%
% 	\\[2ex]
% 	Quindi, per l'unicità data dalla proprietà universale della spanna \((P,p_1,p_2)\) si ottiene \(h\cmp k=\id_P\). Invertendo i ruoli di \(h\) e \(k\), e sfruttando la proprietà universale della spanna \((Q,q_1,q_2)\), si ottiene \(k\cmp h=\id_Q\), e si conclude quindi che la freccia \(h\) è un isomorfismo, con inverso \(h^{-1}=k\).
% \end{proof}
% \begin{remark}[Il prodotto \emph{vs} un prodotto]
% 	Anticipiamo un'osservazione, che riguarda tutti i limiti e tutti i colimiti. Quando ci si riferisce a un prodotto di due oggetti, si utilizza di norma l'articolo indeterminativo (\emph{un} prodotto di \(A\) e \(B\)), poiché la definizione categoriale---cioè mediante una proprietà universale---definisce gli oggetti a meno di un unico isomorfismo. Tuttavia, anche nella letteratura, talvolta si cede alla tentazione di usare l'articolo determinativo (\emph{il} prodotto di \(A\) e \(B\)), dando per scontato il fatto di aver scelto una particolare istanza, ad esempio una specifica costruzione, del prodotto considerato, o più comunemente, riferendosi al concetto generale, come nella frase: ``nella categoria \(\ctSet\), esiste sempre \emph{il} prodotto due insiemi''.  Queste formulazioni, soprattutto la prima, sono da considerarsi un abuso di linguaggio, e perciò andrebbero evitate.
% \end{remark}
% \begin{notation}[Notazione standard del prodotto binario] La notazione standard per il prodotto di due oggetti \(A\) e  \(B\) è la spanna \((A\times B,\pi_A,\pi_B)\):
% 	\[
% 		\begin{tikzcd}
% 			A&A\times B\ar[l, "\pi_A"']\ar[r, "\pi_B"]&B
% 		\end{tikzcd}
% 	\]
% 	Le due frecce uscenti dal prodotto vengono dette \emph{proiezioni}.
% \end{notation}


% \begin{remark}(Sulle proiezioni del prodotto).
% 	Il termine ``proiezioni'' può trarre in inganno, e indurre il lettore a pensare che si tratti, come in \ref{prod_in_Set}, di funzioni suriettive, o più in generale di frecce epi. Questo, in generale, non è vero neanche nella categoria degli insiemi (e delle funzioni). Infatti, si consideri la spanna universale in  \(\ctSet\):
% 	\[
% 		\begin{tikzcd}
% 			A&A\times \emptyset\ar[l, "\pi_A"']\ar[r, "\pi_\emptyset"]&\emptyset
% 		\end{tikzcd}
% 	\]
% 	Poiché ha per codominio l'insieme vuoto, \(\pi_\emptyset\) non può che essere l'identità. Questo forza \(A\times \emptyset=\emptyset\) e dunque \(\pi_A\colon \emptyset\to A\) è suriettiva se, e solo se, \(A=\emptyset\).
% \end{remark}

% \medskip
% Mediante la sua proprietà universale, il prodotto categoriale si estende naturalmente alle frecce. Sia \(\ctC\) una categoria, e \(A,B,C,D\) oggetti tali che esistano i prodotti \(A\times B\) e \(C\times D\). Allora, date due frecce \(f\colon A\to C\) e \(g\colon B\to D\), definiamo \(f\times g\) come l'unica freccia (tratteggiata) che rende commutativo il diagramma seguente:
% \[
% 	\begin{aligned}
% 		\begin{tikzcd}
% 			A\ar[d, "f"']&A\times B\ar[l, "\pi_A"']\ar[r, "\pi_B"]\ar[d, "f\times g", dashed]&B\ar[d, "g"]
% 			\\
% 			C&C\times D\ar[l, "\pi_A"']\ar[r, "\pi_2"]&D
% 		\end{tikzcd}
% 	\end{aligned}
% \]

% In \(\ctSet\), ad esempio, il prodotto di due funzioni \(f,g\) fa esattamente quello che ci si aspetta: per ogni \((a,b)\in A\times B\) si ha che \((f\times g)(a,b)=(f(a),g(b))\).

% \begin{definition}\label{def_cat_con_prodotti}
% 	Una categoria \(\ctC\) ammette (o possiede, o ha) prodotti binari se per ogni coppia di suoi oggetti \(A,B\), esiste una spanna universale sopra di essi.
% \end{definition}

% \subsection{Esempi e non-esempi di prodotti}
% \begin{example}[Monoidi, gruppi, gruppi abeliani]
% 	Le categorie \(\ctMon\), \(\ctCat\) e \(\ctAb\) definite in \ref{ex_cat_monoidi} hanno prodotti. La costruzione canonica del prodotto di due monoidi \(A=(A, \cdot_A, e_A)\) e \(B=(B, \cdot_B, e_B)\)   si ottiene considerando il prodotto Cartesiano \((A\times B,\pi_A, \pi_2)\), ovvero calcolato in \(\ctSet\)), degli insiemi supporto. Esso viene dotato della struttura di monoide ereditata da quelle di \(A\) e \(B\) componente per componente:
% 	\[
% 		(a_1,b_1)\cdot_{A \times B} (a_2,b_2)=(a_1\cdot_A b_1, a_2\cdot_B b_2)\qquad e_{A\times B}=(e_A,e_B)
% 	\]
% 	Le proiezioni risultano essere degli omomorfismi di monoidi.
% 	La stessa cosa vale per gruppi e gruppi abeliani, dove l'inverso si calcola per componenti:
% 	\[
% 		(a,b)^{-1}=(a^{-1},b^{-1})
% 	\]
% 	In modo del tutto analogo si definiscono i prodotti binari nelle categorie di modelli  \ref{ex_cat_sigma_strutture}, come \(\ctRing\), \(\ctRng\), \(\ctcRing\), \(\ctVect\), \dots
% \end{example}
% \color{black}


% \begin{esercizi}
% 	\item
% 	\item
% 	\item
% 	\item
% 	\item
% \end{esercizi}
